\documentclass[11pt,a4paper]{article}

\usepackage{fontspec}
\usepackage{amssymb}
\usepackage{needspace}
\usepackage[spanish,es-noquoting]{babel}
\usepackage[margin=2.6cm]{geometry}
\usepackage{booktabs}
\usepackage{longtable}
\usepackage{enumitem}
\usepackage{xcolor}
\usepackage[colorlinks=true,linkcolor=black,urlcolor=azul]{hyperref}
\usepackage{titlesec}

\definecolor{azul}{HTML}{1F4E79}
\definecolor{gris}{HTML}{5A5A5A}

\titleformat{\section}{\large\bfseries\color{azul}}{\thesection.}{0.5em}{}
\titlespacing{\section}{0pt}{1.4em}{0.6em}
\setlist{itemsep=2pt, topsep=4pt}
\setlength{\parskip}{0.5em}
\setlength{\parindent}{0pt}

\newcommand{\pendiente}{\textcolor{gris}{$\square$}}
\newcommand{\hecho}{\textcolor{azul}{$\blacksquare$}}

\title{\vspace{-1.5cm}\textbf{Canon minero y conflictividad social en el Perú}\\[0.3em]
\large\normalfont Estado del proyecto}
\author{}
\date{Agosto de 2026}

\begin{document}
\maketitle
\vspace{-2em}

\hrule
\vspace{1em}

\section{La pregunta}

\textbf{En los departamentos mineros del Perú, ¿la renta minera reduce la
conflictividad social, o solo cambia aquello por lo que se protesta?}

La pregunta tiene dos mitades que se responden con evidencia distinta. La
primera —\emph{¿hay menos conflicto donde hay más renta?}— se contesta contando:
producción minera, empleo, número de conflictos por departamento y por mes. La
segunda —\emph{¿de qué se pelea la gente?}— no se puede contestar contando. Hay
que leer.

Esa es la apuesta del proyecto: que el interés no está en cuántos conflictos hay,
sino en qué cambia el contenido de las demandas según cuánta renta minera llega
al territorio.

\section{Por qué importa}

La literatura peruana está genuinamente dividida, y las cuatro disciplinas
sociales llegaron a conclusiones distintas mirando los mismos casos.

\begin{itemize}
  \item Los trabajos econométricos sobre Yanacocha encuentran \textbf{efectos
  locales positivos} de la actividad minera sobre el bienestar; otros, mirando un
  horizonte más largo, encuentran que esos efectos \textbf{se desvanecen}.
  \item Los estudios sobre el canon sostienen lo contrario: que descentralizar
  rentas hacia municipios sin capacidad de gestión \textbf{multiplica las
  disputas} por el reparto.
  \item La antropología y la ecología política vienen diciendo, desde hace años,
  que el conflicto \textbf{no es distributivo}: se trata de agua, de territorio y
  de relaciones con el paisaje que no se compensan con dinero.
\end{itemize}

Si la tercera posición es correcta, entonces medir solo transferencias y contar
protestas es mirar el problema con el instrumento equivocado. Este proyecto
intenta poner ambas evidencias en la misma mesa.

\section{Los datos}

Todo está descargado en \texttt{datos/}. Ninguna fuente requiere credenciales.

\Needspace{12\baselineskip}
\begin{longtable}{@{}p{5.1cm}p{6.4cm}p{3.3cm}@{}}
\toprule
\textbf{Fuente} & \textbf{Qué aporta} & \textbf{Formato} \\
\midrule
\endhead
Defensoría del Pueblo &
Reportes mensuales de conflictos: tipo, ubicación, actores y \textbf{descripción
narrativa} de cada caso &
10 PDF (ago 2025 -- jul 2026) \\
\addlinespace
MINEM &
Producción minera y empleo minero por departamento &
ZIP con hojas de cálculo, XLSX \\
\addlinespace
PCM --- Secretaría de Gestión Social &
Mesas de diálogo y 506 compromisos del Estado, con su estado de cumplimiento &
CSV con distrito y coordenadas \\
\addlinespace
Congreso de la República &
251 proyectos de ley sobre minería y canon, con fecha y estado &
JSON \\
\bottomrule
\end{longtable}

El corpus central es el de la Defensoría, porque es el único que dice
\emph{de qué} trata cada conflicto. Lo demás sirve para situarlo.

\section{Qué se ha avanzado}

\begin{itemize}[label=\hecho]
  \item \textbf{Fuentes identificadas y descargadas.} Todas verificadas: se
  comprobó que los archivos abren y contienen lo que dicen contener.
  \item \textbf{Reproducibilidad resuelta.} El script
  \texttt{descargar-datos.sh} rehace la descarga completa y documenta las
  particularidades de cada portal.
  \item \textbf{Bibliografía preliminar reunida} en
  \texttt{bibliografia-borrador.md}, con las referencias centrales de las cuatro
  disciplinas.
\end{itemize}

\section{Qué falta}

\begin{itemize}[label=\pendiente]
  \item \textbf{Convertir los reportes a una base de datos.} Son PDF. No existe
  ninguna versión tabular de esta serie, ni oficial ni de terceros. Hay que
  extraer los casos, uno por uno, a un archivo con el que se pueda trabajar.
  \item \textbf{Clasificar las demandas.} Categorías de partida: agua, tierra y
  territorio, empleo local, renta y compensación, consulta previa,
  incumplimiento de acuerdos anteriores. Habrá que ajustarlas leyendo.
  \item \textbf{Medir el error de esa clasificación.} Ninguna codificación
  —hecha a mano o asistida— sirve si no se sabe cuánto se equivoca. Hay que
  revisar una muestra y contrastar los totales contra las cifras que los propios
  reportes publican.
  \item \textbf{Revisar la bibliografía.} Está armada a las apuradas, copiando de
  buscadores. Hay al menos un año que no cuadra con el volumen citado. Antes de
  que entre a cualquier texto, hay que verificar cada referencia contra una
  fuente primaria.
  \item \textbf{El análisis.} Todavía no hay ninguno.
\end{itemize}

\section{Una decisión metodológica que hay que declarar}

La variable ideal sería el \textbf{canon minero transferido por distrito}. El
Ministerio de Economía lo publica, pero en una aplicación web antigua que no
permite exportar: hay que consultarla a mano, pantalla por pantalla.

Se optó por usar \textbf{producción y empleo minero como aproximación}. Es
defendible, porque el canon corresponde a la mitad del impuesto a la renta que
pagan las empresas mineras y se distribuye a las zonas productoras. Pero es
imperfecto: los precios internacionales introducen ruido, y entre la producción y
la transferencia efectiva hay un rezago de un año o más.

Esto no es un detalle a esconder en una nota al pie. La investigación real está
llena de sustituciones como esta, y lo que distingue a un trabajo serio no es
evitarlas sino \textbf{declararlas y discutir sus límites}.

\vspace{1.5em}
\hrule
\vspace{0.8em}

{\small\color{gris}
\textbf{Cómo está organizada la carpeta.}
\texttt{datos/} contiene las fuentes tal como se descargaron, sin modificar.
\texttt{papers/} reúne la literatura acumulada, con los nombres de archivo con
que llegó cada PDF. \texttt{notas/} son apuntes de trabajo.
\texttt{bibliografia-borrador.md} es la bibliografía pendiente de revisión.

\textbf{Una advertencia sobre las fuentes.} Los portales del Estado peruano
cambian de estructura sin aviso y varios rechazan las descargas automatizadas.
Las particularidades encontradas están anotadas como comentarios dentro de
\texttt{descargar-datos.sh}.
}

\end{document}
